\documentclass[11pt]{article}
\usepackage[T1]{fontenc}
\usepackage{amsmath,amssymb,amsthm}
\usepackage[margin=0.9in]{geometry}
\usepackage{booktabs,array,longtable}
\usepackage[hidelinks]{hyperref}
\newtheorem{theorem}{Theorem}
\newtheorem{corollary}[theorem]{Corollary}
\newcommand{\coeff}[2]{[x^{#1}]\,#2}
\title{Positive expansions for Bala's trinomial coefficient conjecture\
and alternative proofs of the Bajaj--Burns obstructions}
\author{Haobo Ma\\\small ChronoAI Pte Ltd, Singapore
\and Wenlin Zhang\\\small National University of Singapore\\\small ChronoAI Pte Ltd, Singapore}
\date{15 November 2026}
\begin{document}
\maketitle
\begin{abstract}
A finite binomial expansion proves Peter Bala's finite-sum conjecture
(OEIS A351858, FORMULA, revision 2, 27 February 2022; retained in revision 31,
28 October 2025), extends to his rational-function family, and gives alternative
proofs of the family obstructions of Ravi Bajaj and Alexander Burns
(12 August 2026, Sections 2--3).
We obtain a two-parameter coefficient identity and a positive finite formula
for every member of the family, including A351859. We also derive the two
entries' alternating sums, reversion formulas, algebraicity and ordinary Gauss
congruences. The two infinite counterexample families have exact valuation
two, distinguishing these ordinary congruences from the proposed cubic ones.
\end{abstract}

\section{A positive coefficient expansion}
All series are formal. For $H\in\mathbb Q[[x]]$, $[x^n]H$ denotes the
coefficient of $x^n$. A series with constant term one has a multiplicative
inverse. Thus negative integer powers used below are well defined.
For nonnegative integers $a$, $\binom ab=0$ if $b<0$ or $b>a$.

Peter Bala defines
\[
 a(n)=[x^n]\frac{(1+x+x^2)^{3n}}{(1+x)^{2n}},\qquad n\geq0,
\]
and records the conjecture
$a(n)=\sum_{j=0}^{\lfloor n/2\rfloor}\binom{3n}{j}\binom{n-j}{j}$
\cite[FORMULA]{bala2}. The following identity proves it by allowing the
numerator exponent to vary independently of the coefficient index.
The alternating formulas already in Bala's entries arise from expanding
different factors; Section~\ref{sec:siblings} relates them explicitly.
The separate congruence results of Bajaj and Burns
\cite{bbfirst,bb} concern the same family, but not this coefficient identity.

\begin{theorem}\label{thm:main}
For all nonnegative integers $m,n$,
\begin{equation}\label{eq:main}
 [x^n](1+x+x^2)^m(1+x)^{n-m}
 =\sum_{j=0}^{\lfloor n/2\rfloor}\binom mj\binom{n-j}{j}.
\end{equation}
\end{theorem}
\begin{proof}
Use $1+x+x^2=(1+x)(1+x^2/(1+x))$ to obtain
\begin{equation}\label{eq:expand}
 (1+x+x^2)^m(1+x)^{n-m}
 =\sum_{j=0}^{m}\binom mj x^{2j}(1+x)^{n-j}.
\end{equation}
If $2j>n$, the coefficient of $x^n$ in that summand is zero, even when
$n-j<0$, since $(1+x)^{n-j}$ has no negative degrees. Otherwise $n-j\geq0$
and $[x^{n-2j}](1+x)^{n-j}=\binom{n-j}{j}$.
Our binomial convention extends the sum beyond $m$ when necessary.
For $n=0$, both sides equal one.
\end{proof}

\begin{corollary}\label{cor:bala}
Bala's finite-sum conjecture holds for every $n\geq0$:
\[
 a(n)=\sum_{j=0}^{\lfloor n/2\rfloor}\binom{3n}{j}\binom{n-j}{j}.
\]
In particular, $a(n)$ is a positive integer, and its first six terms are
$1,1,7,19,103,376$.
\end{corollary}
\begin{proof}
Set $m=3n$ in Theorem~\ref{thm:main}. All summands are nonnegative and the
$j=0$ summand equals one. Substitution gives the displayed terms.
\end{proof}

For $k\geq1$, put $S_{k-1}(x)=1+x+\cdots+x^{k-1}$ and
\[
 G_k(x)=\frac{S_k(x)^{k+1}}{S_{k-1}(x)^k},\qquad
 u_k(n)=[x^n]G_k(x)^n\quad(n\geq0).
\]
Thus $u_2=a$, and $u_3$ is A351859 \cite{bala3}.

\begin{corollary}\label{cor:family}
For every $k\geq1$ and $n\geq0$,
\begin{equation}\label{eq:family}
 u_k(n)=\sum_{j=0}^{\lfloor n/k\rfloor}\binom{(k+1)n}{j}
 [x^{n-kj}]S_{k-1}(x)^{n-j}.
\end{equation}
Every $u_k(n)$ is a positive integer. In particular,
$u_1(n)=\binom{2n}{n}$ and
\[
 u_3(n)=\sum_{j=0}^{\lfloor n/3\rfloor}\binom{4n}{j}
 [x^{n-3j}](1+x+x^2)^{n-j}.
\]
\end{corollary}
\begin{proof}
The identity $S_k=S_{k-1}+x^k$ gives
\[
 G_k(x)^n=\sum_{j=0}^{(k+1)n}\binom{(k+1)n}{j}
 x^{kj}S_{k-1}(x)^{n-j}.
\]
Terms with $kj>n$ cannot contribute. For every other term,
$j\leq n/k\leq n$, so its remaining factor is a polynomial with
nonnegative integer coefficients. If $k\geq2$, the $j=0$ term is positive:
choosing $x$ from each factor of $S_{k-1}^n$ contributes one. For $k=1$,
$S_0=1$ and only $j=n$ contributes, giving $\binom{2n}{n}$. This also
handles $n=0$.
\end{proof}

The cubic formula is entirely explicit: for $0\leq d\leq2N$,
\[
 [x^d](1+x+x^2)^N
 =\sum_{h=\max(0,d-N)}^{\lfloor d/2\rfloor}
 \frac{N!}{h!(d-2h)!(N-d+h)!}.
\]
Here $h$ factors contribute $x^2$, $d-2h$ contribute $x$, and the rest
contribute one. Outside this range the coefficient is zero.

\section{Other formulas in the two entries}\label{sec:siblings}
The following corollaries supply direct proofs of formulas already recorded
in A351858 and A351859. Their attribution remains with the source entries.

\begin{corollary}\label{cor:alternating}
For $n\geq1$, the alternating sums in the two entries are
\begin{align}
 u_2(n)&=\sum_{h=0}^{n}\sum_{j=0}^{n-h}
 (-1)^{n-h-j}\binom nh\binom{3n}{j}
 \binom{4n-2j-h-1}{n-h-j},\label{eq:double}\\
 u_3(n)&=\sum_{i=0}^{n}\sum_{j=0}^{n}\sum_{h=0}^{j}
 (-1)^j\binom{4n}{n-2i-j-h}\binom{4n}{i}
 \binom{3n+j-1}{j}\binom jh.\label{eq:triple}
\end{align}
At $n=0$, both sequences have value one.
\end{corollary}
\begin{proof}
For \eqref{eq:double}, write
\[
 G_2(x)^n=(1+x)^n\bigl(1+x(1+x)\bigr)^{3n}(1+x)^{-3n}.
\]
Choose degrees $h$ and $j$ in the first two binomial expansions.
The remaining factor is $(1+x)^{j-3n}$, whose degree $n-h-j$
coefficient is
$(-1)^{n-h-j}\binom{4n-2j-h-1}{n-h-j}$.
For \eqref{eq:triple}, use
\[
 G_3(x)^n=(1+x)^{4n}(1+x^2)^{4n}
 \sum_{j\geq0}(-1)^j\binom{3n+j-1}{j}x^j(1+x)^j.
\]
Expanding the second factor at degree $2i$ and the last at degree $h$
leaves degree $n-2i-j-h$ in $(1+x)^{4n}$. Terms with $i>n$ or $j>n$
cannot contribute. This proves the stated finite ranges.
\end{proof}

\begin{corollary}\label{cor:reversion}
Let $T_k(x)$ be the compositional inverse of $x/G_k(x)$, let
$F_k(x)=T_k(x)/x$, and set $A_k(x)=\sum_{n\geq0}u_k(n)x^n$.
Then
\begin{equation}\label{eq:reversion}
 A_k=1+x\frac{F_k'}{F_k},\qquad
 F_k S_{k-1}(xF_k)^k=S_k(xF_k)^{k+1}.
\end{equation}
Thus $A_k$ is algebraic over $\mathbb Q(x)$, and
\[
 A_k(z)=\operatorname{Diag}\frac{1}{1-tG_k(x)}
 :=\sum_{n\geq0}[x^nt^n]\frac{z^n}{1-tG_k(x)}.
\]
These statements include the reversion and diagonal formulas in both entries.
\end{corollary}
\begin{proof}
More generally let $\phi\in1+x\mathbb Q[[x]]$, let $T=x\phi(T)$ and
$F=T/x=\phi(T)$. Write $\operatorname{Res}_t$ for the coefficient of
$t^{-1}$ in a Laurent series multiplied by $dt$.
For $n\geq1$, integration by parts and the substitution $x=t/\phi(t)$ give
\begin{align*}
 [x^n]\log F
 &=\frac1n\operatorname{Res}_x x^{-n}\,d\log F\\
 &=\frac1n\operatorname{Res}_t
 \frac{\phi(t)^{n-1}\phi'(t)}{t^n}\,dt
 =\frac1n[t^n]\phi(t)^n.
\end{align*}
In the last step use $(\phi^n)'=n\phi^{n-1}\phi'$.
Formal residues are unchanged by a substitution with nonzero linear term:
this follows for monomials from the residue of a derivative being zero,
and for $t^{-1}dt$ from the logarithmic derivative having residue one.
Differentiating the coefficient identity proves the first equation in
\eqref{eq:reversion}. For $\phi=G_k$, the equation $F_k=G_k(xF_k)$
gives the second. It is a nonzero polynomial equation (at $x=0$ it is
$F_k-1=0$). Hence $F_k$, its derivative, and $A_k$ are algebraic in
characteristic zero. Finally, expand $(1-tG_k(x))^{-1}=\sum_{n\geq0}t^nG_k(x)^n$
to obtain the diagonal identity.
\end{proof}

\begin{corollary}\label{cor:gauss}
Let $H\in1+x\mathbb Z[[x]]$, let $c\geq1$ be an integer and put
$v(n)=[x^{cn}]H(x)^n$. For every prime $p$ and integers $n,r\geq1$,
\[
 v(np^r)\equiv v(np^{r-1})\pmod{p^r}.
\]
In particular this holds for
$v(n)=[x^{cn}]G_k(x)^{sn}$ for all $k,c\geq1$ and $s\in\mathbb Z$.
\end{corollary}
\begin{proof}
Coefficientwise Frobenius gives $H(x)^p\equiv H(x^p)\pmod p$.
If $U\equiv V\pmod{p^e}$ with $e\geq1$, the binomial expansion gives
$U^p\equiv V^p\pmod{p^{e+1}}$: the linear term has the required factor,
the intermediate binomial coefficients are divisible by $p$, and
$ep\geq e+1$ for the final term.
Iteration and then taking an $n$th power yield
$H(x)^{np^r}\equiv H(x^p)^{np^{r-1}}\pmod{p^r}$.
Extract degree $cnp^r$. Finally $G_k^s$ has integer coefficients for
every integer $s$, since $S_{k-1}$ and $S_k$ are integral units.
\end{proof}

\begin{corollary}\label{cor:criterion}
For any integer sequence $(g(n))_{n\geq1}$, there is a unique
$\phi\in1+x\mathbb Q[[x]]$ such that $g(n)=[x^n]\phi(x)^n$.
This $\phi$ has integer coefficients if and only if $g$ satisfies the
ordinary Gauss congruences for every prime and all positive indices.
\end{corollary}
\begin{proof}
Set $F(x)=\exp(\sum_{n\geq1}g(n)x^n/n)$ and $T=xF$.
The inverse of $T$ exists, and $\phi=F\circ T^{-1}$ satisfies $T=x\phi(T)$.
The coefficient calculation in Corollary~\ref{cor:reversion} then gives
$[x^n]\phi^n=n[x^n]\log F=g(n)$; it also proves uniqueness.
If $\phi$ is integral, Corollary~\ref{cor:gauss} gives the congruences.

Conversely assume those congruences and define
$b_d=d^{-1}\sum_{e\mid d}\mu(d/e)g(e)$, where $\mu$ is the M\"obius
function (zero on integers divisible by a prime square, and $(-1)^j$ on
products of $j$ distinct primes).
For $d=p^a m$ with $p\nmid m$, the numerator equals
\[
 \sum_{e\mid m}\mu(m/e)\bigl(g(ep^a)-g(ep^{a-1})\bigr),
\]
which is divisible by $p^a$. Doing this for every prime divisor of $d$
shows $b_d\in\mathbb Z$. M\"obius inversion gives
$\sum_{d\mid n}d b_d=g(n)$, so taking formal logarithms proves
\[
 F(x)=\prod_{d\geq1}(1-x^d)^{-b_d}\in1+x\mathbb Z[[x]].
\]
The product is well defined degree by degree. The inverse of $T=xF$
has integer coefficients by recursive coefficient comparison. Therefore
$\phi=F\circ T^{-1}$ is integral, as required.
\end{proof}

\section{The cubic-congruence obstruction}
The COMMENTS of A351858 assert: ``Conjecture: for $k\geq2$, each sequence
$u_k$ satisfies the same supercongruences as the central binomial
coefficients.'' Explicitly, these are
\begin{equation}\label{eq:cubic}
 u_k(np^r)\equiv u_k(np^{r-1})\pmod{p^{3r}}
 \quad(p\geq5\text{ prime};\ n,r\geq1).
\end{equation}
The following obstructions were given by Bajaj and Burns
\cite[Sections 2--3]{bbfirst}; their subsequent version \cite{bb} adds
a separate proof of the fixed $k=2$ cubic congruences in Sections 5--7.
Here the positive expansion gives an alternative derivation of the obstructions.

\begin{corollary}[Bajaj--Burns]\label{cor:obstruction}
For every prime $p\geq5$,
\[
 u_{p-1}(p)=\binom{2p-1}{p}+p^2(p-2),\qquad
 u_p(p)=\binom{2p-1}{p}+p^2.
\]
Since $u_{p-1}(1)=u_p(1)=1$, both differences from the index-one value
have $p$-adic valuation exactly two. Thus \eqref{eq:cubic} fails for
these parameter choices.
\end{corollary}
\begin{proof}
For $k=p-1$, \eqref{eq:family} has terms $j=0,1$. The first counts
weak compositions of $p$ into $p$ parts with each part at most $p-2$:
\[
 [x^p]S_{p-2}(x)^p=\binom{2p-1}{p}-p^2.
\]
Indeed a specified part at least $p-1$ leaves one to distribute into
$p$ parts, and two such large parts cannot occur. The second term is
$p^2[x]S_{p-2}^{p-1}=p^2(p-1)$. For $k=p$, the first term is
$\binom{2p-1}{p}-p$, while the second is $p(p+1)$.

For completeness, the elementary Wolstenholme congruence needed here is
$\binom{2p-1}{p}\equiv1\pmod{p^3}$. In the localization of $\mathbb Z$
at $p$, put $H_i=\sum_{a=1}^{p-1}a^{-i}$.
Inversion permutes the nonzero residues, so $H_2\equiv\sum_{a=1}^{p-1}a^2
\equiv0\pmod p$ for $p\geq5$.
Pairing $a$ and $p-a$ gives $2H_1\equiv-pH_2\equiv0\pmod{p^2}$.
Consequently
\[
 \binom{2p-1}{p}=\prod_{a=1}^{p-1}(1+p/a)
 \equiv1+pH_1+\frac{p^2}{2}(H_1^2-H_2)\equiv1\pmod{p^3}.
\]
The two differences are therefore congruent to $p^2(p-2)$ and $p^2$
modulo $p^3$, respectively, with nonzero quotients modulo $p$.
\end{proof}

At $p=5$ this recovers $u_4(5)=201$ and $u_5(5)=151$, with differences
$200$ and $150$. The fixed $k=2$ and $k=3$ statements are not contradicted:
neither can equal $p-1$ or $p$ for $p\geq5$.
The further coefficient-slope sentence in both OEIS entries says
``may satisfy the same supercongruences.'' It does not assert a universal
theorem. A universal reading across all $k,c,s$ fails at $c=s=1$, $k=4,5$;
this is not a refutation of an assertion made in that sentence, nor does it
decide the fixed $k=3$ suggestion in A351859.

\section{Coverage and verification}
\small
\begin{longtable}{@{}>{\raggedright\arraybackslash}p{0.46\linewidth}>{\raggedright\arraybackslash}p{0.48\linewidth}@{}}
\toprule
Source item & Coverage\\
\midrule
\endhead
A351858, FORMULA finite-sum conjecture & Proved, Corollary~\ref{cor:bala}; all $n\geq0$.\\
A351858 COMMENTS family; A351859 NAME & Positive formula and integrality proved, Corollary~\ref{cor:family}; a generalization, not separately posed as a conjecture.\\
Alternating sums, both FORMULA fields & Reproved; prior formulas Bala (2022), Corollary~\ref{cor:alternating}.\\
Reversion and algebraicity, both FORMULA fields & Reproved; prior formulas Bala (2022), Corollary~\ref{cor:reversion}.\\
Ordinary Gauss congruences, both COMMENTS fields & Reproved; prior statements Bala (2022), Corollary~\ref{cor:gauss}, also for all integral slopes.\\
General rational-series existence and integral-coefficient criterion, A351858 COMMENTS & Reproved; prior statement Bala (2022), Corollary~\ref{cor:criterion}.\\
All-$k$ cubic conjecture, A351858 COMMENTS & Refuted (explicit counterexamples); reproved, prior proof Bajaj--Burns (2026), Corollary~\ref{cor:obstruction}. No corrected cubic theorem is asserted here.\\
Fixed $k=2$ cubic congruences & Prior proof cited: Bajaj--Burns (2026), Sections 5--7 of \cite{bb}; separate from our coefficient proof.\\
Fixed $k=3$ cubic congruences, A351859 COMMENTS & Not covered by these proofs.\\
Coefficient-slope ``may satisfy'' suggestions & A universal all-$k$ assertion is not asserted in the source; fails for $k=4,5$, $c=s=1$, $p=5$. The fixed $k=3$ suggestion is not covered.\\
A351858 asymptotic formula (Kotesovec, 2022) & Not covered.\\
Cross-referenced ordinary Gauss statements & Corollary~\ref{cor:gauss} applies to the coefficient representations in A000984, A103885, A333090, A333564, A333579, A333715, A351856, A351857 and A372211--A372213. These are prior statements, not new conjecture resolutions.\\
Other cross-reference items & Cubic conjectures, recurrences and zero-location conjectures in those entries, and A333565's other formulas, are not covered.\\
\bottomrule
\end{longtable}
\normalsize
The cross-reference numbers in the table denote the corresponding
\href{https://oeis.org}{OEIS} entries. For A333090 the integral series is
$(1+x)S(x/(1+x))$, where $S$ is the large Schr\"oder series; for A333564
it is the quotient of the Catalan series at $x$ and $-x$. Both are integral
units. The remaining listed representations are rational integral units,
with the slope and power specified by their NAME or FORMULA fields.
The table does not assign a common status to their distinct cubic conjectures.

The Lean declaration
\begin{quote}\small
\texttt{D5.S1.Recurrence.}\\
\texttt{BalaTrinomialCoefficientIdentity.result}
\end{quote}
in the \href{https://github.com/the-omega-institute/trureturing/blob/4a7117cee792c7d5598d7f4770577d557fb45137/D5/S1/Recurrence/BalaTrinomialCoefficientIdentity.lean}{pinned Lean file}
has exactly the statement of Corollary~\ref{cor:bala}, for every natural
$n$, including zero. Its definition \texttt{source} is the rational quotient,
formed independently of the sum; the natural-number sum is cast into
$\mathbb Q$. The two-parameter identity, positive family expansion,
other-formula corollaries and alternative obstruction proofs have the written
proofs above. The accompanying \texttt{verify.py} supplies finite exact-integer
coefficient comparisons, whose ranges and results are in \texttt{verification.json}.

\paragraph{Availability and license.}
The source package contains this \LaTeX{} source, an exact-integer checker,
coefficient reference values, metadata and a rebuild guide.
This note is licensed under Creative Commons Attribution 4.0 International
(CC BY 4.0).

\begin{thebibliography}{9}
\bibitem{bala2}
Peter Bala, \emph{A351858}, The On-Line Encyclopedia of Integer Sequences,
submitted 27 February 2022; FORMULA conjecture added in revision 2;
revision 31 (28 October 2025), NAME, COMMENTS and FORMULA.
\url{https://oeis.org/A351858};
\href{https://oeis.org/history?seq=A351858\&start=20}{original edit history}.

\bibitem{bala3}
Peter Bala, \emph{A351859}, The On-Line Encyclopedia of Integer Sequences,
submitted 1 March 2022, revision 24 (30 May 2025), NAME, COMMENTS and FORMULA.
\url{https://oeis.org/A351859}.

\bibitem{bbfirst}
Ravi Bajaj and Alexander Burns,
\emph{An exact boundary for the A351858 cyclotomic family},
12 August 2026, Sections 2--3,
\href{https://github.com/rbajaj5/a183068-supercongruence/blob/524c02ec2a8da45b9fc7078e7de6f3a044360e82/related-results/A351858CyclotomicFamilyBoundary.md}{initial version}.

\bibitem{bb}
Ravi Bajaj and Alexander Burns,
\emph{The A351858 cubic tower and its exact family boundary},
12 August 2026, Sections 5--7,
\href{https://github.com/rbajaj5/a183068-supercongruence/blob/dd9762d8cdf9718ecc09637dc470e87100677886/related-results/A351858CyclotomicFamilyBoundary.md}{subsequent version}.
\end{thebibliography}
\end{document}
